\documentclass[../../main.tex]{subfiles}

\begin{document}

\section{连通空间}\label{section:ch03-connected-spaces}

拓扑空间中连通性的定义是易于理解的。如果空间能够拆成两个``团''——两个无交的开集，我们就说空间是可以``分割''的。否则，便称它是连通的。这样一个简单的想法引发了后续的相关讨论。

\begin{definition}\label{definition:ch03-separation}
设 $X$ 是一个拓扑空间。所谓 $X$ 的一个\textbf{分割}（separation），是指 $X$ 的一对无交的非空开子集 $U$ 和 $V$，它们的并等于 $X$。如果 $X$ 的分割不存在，则称空间 $X$ 是\textbf{连通的}（connected）。
\end{definition}

连通性显然是一个拓扑性质，因为它的定义仅涉及 $X$ 的开集族。换句话说，如果 $X$ 是连通的，那么与 $X$ 同胚的每一空间都是连通的。

连通性的定义也可以用以下方式给出：空间 $X$ 是连通的当且仅当 $X$ 中既开又闭的子集只有空集和 $X$ 自身。事实上，若 $A$ 是 $X$ 中一个既开又闭的非空真子集，那么，$U=A$ 和 $V=X\setminus A$ 是 $X$ 中的非空无交开集使得其并等于 $X$，从而它们构成了 $X$ 的一个分割。反之，如果 $U$ 和 $V$ 构成 $X$ 的一个分割，则 $U$ 便是 $X$ 的既开又闭的非空真子集。

对于拓扑空间 $X$ 的子空间 $Y$，还有另一种定义其连通性的有用方法：

\begin{lemma}\label{lemma:ch03-subspace-separation}
如果 $Y$ 是 $X$ 的子空间，则 $Y$ 的一个分割是一对无交的非空集合 $A$ 和 $B$，它们的并等于 $Y$，并且 $A$ 和 $B$ 中的任何一个都不包含另一个的极限点。如果空间 $Y$ 不存在这样的分割，则空间 $Y$ 是连通的。
\end{lemma}

\begin{proof}
首先，设 $A$ 和 $B$ 是 $Y$ 的一个分割，则 $A$ 在 $Y$ 中既开又闭。$A$ 在 $Y$ 中的闭包是 $\overline{A}\cap Y$（按照惯例，这里 $\overline{A}$ 表示 $A$ 在 $X$ 中的闭包）。由于 $A$ 在 $Y$ 中是闭的，所以 $A=\overline{A}\cap Y$，因此 $\overline{A}\cap B=\varnothing$。由于 $\overline{A}$ 等于 $A$ 与其极限点集的并，所以 $B$ 不包含 $A$ 的极限点。同理，$A$ 也不包含 $B$ 的极限点。

反之，假设 $A$ 和 $B$ 是 $Y$ 中两个非空无交集合，其并等于 $Y$，并且其中任何一个不包含另一个的极限点。于是 $\overline{A}\cap B=\varnothing$，$A\cap\overline{B}=\varnothing$。由此可见 $\overline{A}\cap Y=A$ 和 $\overline{B}\cap Y=B$，于是 $A$ 和 $B$ 都是 $Y$ 中的闭集，再根据 $A=Y\setminus B$ 和 $B=Y\setminus A$，便可见它们也都是 $Y$ 中的开集。
\end{proof}

\begin{example}\label{example:ch03-two-point-indiscrete}
设 $X$ 是只有两个点的密着拓扑空间。显然，$X$ 的分割不存在，因此 $X$ 是连通的。
\end{example}

\begin{example}\label{example:ch03-minus-two-points}
设 $Y$ 是实直线 $\mathbb{R}$ 的子空间 $[-1,0)\cup(0,1]$。由于 $[-1,0)$ 与 $(0,1]$ 都是 $Y$ 中非空的开集（尽管它们不是 $\mathbb{R}$ 中的开集），因此，它们构成 $Y$ 的一个分割。另外请注意，它们中任何一个集合都不包含另一个集合的极限点。（$0$ 是它们的一个公共极限点，但这个点不在子空间 $Y$ 中。）
\end{example}

\begin{example}\label{example:ch03-closed-interval}
设 $X$ 是实直线 $\mathbb{R}$ 的子空间 $[-1,1]$。$[-1,0]$ 和 $(0,1]$ 是两个无交非空集合，但它们不构成 $X$ 的分割，因为第一个集合不是 $X$ 中的开集。另外请注意，第一个集合包含了第二个集合的极限点 $0$。事实上，空间 $[-1,1]$ 的分割是不存在的。稍后我们将证明这一点。
\end{example}

\begin{example}\label{example:ch03-rationals}
有理数集 $\mathbb{Q}$ 是不连通的。事实上，$\mathbb{Q}$ 的连通子空间只有单点集。因为如果 $Y$ 是 $\mathbb{Q}$ 中包含点 $p$ 和 $q$ 的一个子空间，则在 $p$ 和 $q$ 之间可以选择一个无理数 $a$，并且 $Y$ 可以表示为两个开集
\begin{align*}
Y\cap(-\infty,a)\quad\text{和}\quad Y\cap(a,+\infty)
\end{align*}
的并。
\end{example}

\begin{example}\label{example:ch03-graph-asymptote}
考虑平面 $\mathbb{R}^{2}$ 的子集：
\begin{align*}
X=\{x\times y\mid y=0\}\cup\{x\times y\mid x>0\text{ 且 }y=1/x\}.
\end{align*}
$X$ 是不连通的。事实上，上面给出的两个集合就是 $X$ 的一个分割，因为它们中的任何一个都不包含另一个的极限点（见\cref{figure:ch03-connected-fig1}）。
\end{example}

我们已经给出了几个不连通空间的例子。然而，怎样构造连通空间呢？我们即将证明几个定理，它们给出了一些由已知连通空间构造新的连通空间的方法。在下一节，我们将应用这些定理证明几个特殊空间的连通性，比如 $\mathbb{R}$ 中的区间，$\mathbb{R}^{n}$ 中的球和方体。首先给出下述引理：

\begin{lemma}\label{lemma:ch03-in-separation}
如果集合 $C$ 与 $D$ 构成 $X$ 的一个分割，并且 $Y$ 是 $X$ 的一个连通子空间，那么，$Y$ 或者包含于 $C$，或者包含于 $D$。
\end{lemma}

\begin{proof}
由于 $C$ 和 $D$ 都是 $X$ 中的开集，从而 $C\cap Y$ 和 $D\cap Y$ 都是 $Y$ 中的开集。这两个集合无交且它们的并等于 $Y$。倘若它们中每一个都不是空集，则它们便构成了 $Y$ 的一个分割。因此，它们之中必有一个为空集。于是，或者 $Y$ 包含于 $C$，或者 $Y$ 包含于 $D$。
\end{proof}

\par
\begin{figure}[htbp]
\centering
\includegraphics[width=0.5\textwidth]{ch03-fig23-1.png}
\caption{例\,\cref{example:ch03-graph-asymptote}中的不连通空间}\label{figure:ch03-connected-fig1}
\end{figure}

\begin{theorem}\label{theorem:ch03-union-common-point}
含一个公共点的 $X$ 的连通子空间族的并是连通的。
\end{theorem}

\begin{proof}
设 $\{A_{\alpha}\}$ 是空间 $X$ 中连通子空间的一个族。$p$ 是 $\bigcap A_{\alpha}$ 的一个点。我们证明空间 $Y=\bigcup A_{\alpha}$ 是连通的。设 $Y=C\cup D$ 是 $Y$ 的一个分割。那么点 $p$ 将属于 $C$ 或 $D$。不妨设 $p\in C$。因为 $A_{\alpha}$ 是连通的，它必然整个地包含于 $C$ 或 $D$，而由于它包含 $C$ 中的点 $p$，所以 $A_{\alpha}$ 不可能包含在 $D$ 中。因此对于每一个 $\alpha$，有 $A_{\alpha}\subset C$，于是 $\bigcup A_{\alpha}\subset C$，这与 $D$ 非空矛盾。
\end{proof}

\begin{theorem}\label{theorem:ch03-closure-connected}
设 $A$ 是 $X$ 的一个连通子空间。若 $A\subset B\subset\overline{A}$，则 $B$ 也是连通的。
\end{theorem}

换句话说，如果 $B$ 等于连通子空间 $A$ 加上它的部分或全部极限点，那么 $B$ 是连通的。

\begin{proof}
设 $A$ 连通且 $A\subset B\subset\overline{A}$。如果 $B=C\cup D$ 是 $B$ 的一个分割。根据\cref{lemma:ch03-in-separation}，$A$ 必定整个地包含于 $C$ 或者包含于 $D$，假设 $A\subset C$。于是 $\overline{A}\subset\overline{C}$。由于 $C$ 与 $D$ 无交，所以 $B$ 与 $D$ 无交，这与 $D$ 是 $B$ 的非空子集矛盾。
\end{proof}

\begin{theorem}\label{theorem:ch03-continuous-image}
连通空间在连续映射下的像是连通的。
\end{theorem}

\begin{proof}
设 $f\colon X\to Y$ 是一个连续映射，$X$ 是连通空间。我们来证明它的像空间 $Z=f(X)$ 是连通的。因为把 $f$ 的值域限制在空间 $Z$ 上而得到的映射也是连续的，所以只要考虑下述连续满映射
\begin{align*}
g\colon X\longrightarrow Z
\end{align*}
就行了。假设 $Z=A\cup B$ 是由 $Z$ 中两个非空无交开集构成的 $Z$ 的一个分割。那么 $g^{-1}(A)$ 与 $g^{-1}(B)$ 也无交，并且其并等于 $X$。由于 $g$ 是一个连续映射，所以它们都是 $X$ 中的开集。由于 $g$ 是满映射，所以它们都是非空的。因此 $g^{-1}(A)$ 与 $g^{-1}(B)$ 便构成了 $X$ 的一个分割，这与 $X$ 的连通性矛盾。
\end{proof}

\begin{theorem}\label{theorem:ch03-finite-product}
有限多个连通空间的笛卡儿积是连通的。
\end{theorem}

\begin{proof}
首先对两个连通空间 $X$ 与 $Y$ 的积加以证明。这个证明可以说得直观些。在积空间 $X\times Y$ 中选取一个``基点''$a\times b$。注意到``水平薄片''$X\times b$ 与 $X$ 同胚，因而是连通的；每一个``竖立薄片''$x\times Y$ 与 $Y$ 同胚，因而也是连通的。于是，每一个``十字型''空间
\begin{align*}
T_{x}=(X\times b)\cup(x\times Y)
\end{align*}
是两个具有公共点 $x\times b$ 的连通空间之并，因而也是连通的。见\cref{figure:ch03-connected-fig2}。现在考虑所有这些十字型空间的并 $\bigcup_{x\in X}T_{x}$。这个并是连通的，因为它是含有公共点 $a\times b$ 的连通空间族的并。由于这个并等于 $X\times Y$，因此空间 $X\times Y$ 连通。

由于（易证）$X_{1}\times\cdots\times X_{n}$ 同胚于 $(X_{1}\times\cdots\times X_{n-1})\times X_{n}$，所以对于有限多个连通空间的笛卡儿积，我们可以用归纳法给出证明。
\end{proof}

\par
\begin{figure}[htbp]
\centering
\includegraphics[width=0.5\textwidth]{ch03-fig23-2.png}
\caption{有限笛卡儿积的连通性（十字型空间）}\label{figure:ch03-connected-fig2}
\end{figure}

人们自然会问：上述定理是否可以推广到任意多个连通空间的笛卡儿积上？下面的例子告诉我们，这与笛卡儿积上所选择的拓扑有关。

\begin{example}\label{example:ch03-box-disconnected}
考虑 $\mathbb{R}^{\omega}$ 上的箱拓扑。$\mathbb{R}^{\omega}$ 可以表示成所有有界实数序列的集合 $A$ 与所有无界序列的集合 $B$ 的并。这两个集合无交，且每一集合都是箱拓扑中的开集。这是由于任意选取 $\mathbb{R}^{\omega}$ 的一个点 $a$，若 $a$ 是一个有界序列，则
\begin{align*}
U=(a_{1}-1,a_{1}+1)\times(a_{2}-1,a_{2}+1)\times\cdots
\end{align*}
是由有界序列组成的一个开集。而当 $a$ 为无界序列时，则上述 $U$ 是由无界序列组成的一个开集。由此可见，尽管 $\mathbb{R}$ 是连通的（下节我们将给出证明），但 $\mathbb{R}^{\omega}$ 关于箱拓扑却不是连通的。
\end{example}

\begin{example}\label{example:ch03-product-connected}
考虑 $\mathbb{R}^{\omega}$ 的积拓扑。假设 $\mathbb{R}$ 是连通的，我们来证明 $\mathbb{R}^{\omega}$ 是连通的。设 $\widetilde{\mathbb{R}}^{n}$ 是由所有序列 $\boldsymbol{x}=(x_{1},x_{2},\cdots)$ 组成的 $\mathbb{R}^{\omega}$ 的子空间，其中，当 $i>n$ 时有 $x_{i}=0$。显然，$\widetilde{\mathbb{R}}^{n}$ 同胚于 $\mathbb{R}^{n}$，根据前一个定理可见，它是连通的。由于每一个 $\widetilde{\mathbb{R}}^{n}$ 都包含点 $\boldsymbol{0}=(0,0,\cdots)$，所以所有这些 $\widetilde{\mathbb{R}}^{n}$ 的并 $\mathbb{R}^{\infty}$ 是连通的。我们来证明 $\mathbb{R}^{\infty}$ 的闭包等于 $\mathbb{R}^{\omega}$，由此推出 $\mathbb{R}^{\omega}$ 也是连通的。

设 $\boldsymbol{a}=(a_{1},a_{2},\cdots)$ 是 $\mathbb{R}^{\omega}$ 的一个点，$U=\prod U_{i}$ 为积拓扑中含有点 $\boldsymbol{a}$ 的一个基元素。我们证明 $U$ 与 $\mathbb{R}^{\infty}$ 有交。选取一个整数 $N$，使得当 $i>N$ 时有 $U_{i}=\mathbb{R}$。则 $\mathbb{R}^{\infty}$ 的点
\begin{align*}
\boldsymbol{x}=(a_{1},\cdots,a_{n},0,0,\cdots)
\end{align*}
属于 $U$，这是由于对于所有的 $i$ 有 $a_{i}\in U_{i}$，并且对于所有的 $i>N$ 有 $0\in U_{i}$。
\end{example}

上述讨论的推广说明在积拓扑下任意多个连通空间的笛卡儿积是连通的。由于我们不需要这个结果，其证明留作习题。

\subsection*{习题}

\begin{enumerate}
\item 设 $\mathscr{T}$ 和 $\mathscr{T}'$ 是 $X$ 的两个拓扑。若 $\mathscr{T}'\subset\mathscr{T}$，试问关于两种拓扑的连通性之间有什么联系？
\item 设 $\{A_{n}\}$ 是 $X$ 的连通子空间的一个序列，并且对于所有的 $n$ 有 $A_{n}\cap A_{n+1}\neq\varnothing$。证明 $\bigcup A_{n}$ 是连通的。
\item 设 $\{A_{\alpha}\}$ 是 $X$ 的连通子空间的一个族。$A$ 是 $X$ 的一个连通子空间。证明：若对于每一个 $\alpha$ 有 $A\cap A_{\alpha}\neq\varnothing$，则 $A\cup(\bigcup A_{\alpha})$ 是连通的。
\item 证明：若 $X$ 为无限集，则 $X$ 关于有限补拓扑是连通的。
\item 若一个空间的连通子空间只有单点集，则这个空间称为\textbf{完全不连通}（totally disconnected）空间。证明：若 $X$ 具有离散拓扑，则 $X$ 是完全不连通的。其逆成立吗？
\item 设 $A\subset X$。证明：若 $C$ 是 $X$ 的一个连通子空间，并且 $C$ 与 $A$ 和 $X\setminus A$ 都有交，则 $C$ 与 $\mathrm{Bd}\,A$ 也有交。
\item 空间 $\mathbb{R}_{\ell}$ 连通吗？验证你的结论。
\item 判定 $\mathbb{R}^{\omega}$ 关于一致拓扑是否连通。
\item 设 $A$ 是 $X$ 的一个真子集，$B$ 是 $Y$ 的一个真子集。若 $X$ 和 $Y$ 都是连通的，证明
\begin{align*}
(X\times Y)\setminus(A\times B)
\end{align*}
是连通的。
\item 设 $\{X_{\alpha}\}_{\alpha\in J}$ 是连通空间的一个加标族，$X$ 是积空间
\begin{align*}
X=\prod_{\alpha\in J}X_{\alpha}.
\end{align*}
设 $\boldsymbol{a}=(a_{\alpha})$ 是 $X$ 的一个给定的点。
\begin{enumerate}
\item 对于 $J$ 的任何一个有限子集 $K$，以 $X_{K}$ 表示由所有那些点 $\boldsymbol{x}=(x_{\alpha})$ 的集合，其中当 $\alpha\notin K$ 时有 $x_{\alpha}=a_{\alpha}$。证明 $X_{K}$ 是连通的。
\item 所有空间 $X_{K}$ 的并 $Y$ 是连通的。
\item 证明：$X$ 等于 $Y$ 的闭包。从而 $X$ 是连通的。
\end{enumerate}
\item 设 $p\colon X\to Y$ 是一个商映射。证明：若每一个 $p^{-1}(\{y\})$ 是连通的，并且 $Y$ 也是连通的，则 $X$ 是连通的。
\item 设 $Y\subset X$，$X$ 和 $Y$ 都是连通的。证明：若 $A$ 和 $B$ 构成 $X\setminus Y$ 的一个分割，则 $Y\cup A$ 和 $Y\cup B$ 都是连通的。
\end{enumerate}

\end{document}